\section{Introduction}

Digital holographic microscopes (DHMs) have been widely applied in material and biological
applications. For instance, among many biological and biomedical applications, DHM
systems are used for analyzing cells and tissues, as well as disease diagnosis and screening
 and assessing the polarimetric properties of biomedical samples.
The performance of DHM technologies relies heavily on computational reconstruction processing to provide trustworthy
sample information. The required computational reconstruction algorithms are uniquely
dependent on the optical configuration of the DHM system. An incorrect selection of the
reconstruction algorithms leads to distorted and inaccurate amplitude and phase measurements. 

DHM systems record the interference pattern (e.g., hologram) generated between the scattered light from the sample, named object wavefront, and a known reference wavefront.
DHM systems operate in off-axis, slightly off-axis, or in-line (also known as on-axis) configurations based on the interference angle.
Therefore, the selection of the DHM reconstruction
algorithms depends on this interference angle. 

For example, off-axis DHM systems enable the
reconstruction of the complex amplitude distribution of an object wavefront from a single hologram since the three components of the hologram are entirely separable in the Fourier
domain.

Therefore, the reconstruction algorithm in off-axis DHM systems requires the spatial filtering of the sample frequencies from the hologram’s spectrum. In contrast, the
spectral components of a hologram in the Fourier domain overlay totally or partially in in-line
and slightly off-axis DHM systems, respectively, requiring the acquisition of multiple phase-shifted holograms and the application of phase-shifting (PS) techniques [1, 26] to reconstruct the desired object information. 

In addition, the selection of the DHM reconstruction algo-
rithms also depends on whether the DHM imaging system operates in telecentric or non-telecentric regimes.
 DHM systems operating in the telecentric regime only require the interference angle compensation in the off-axis and slightly off-axis configuration.
Oppositely, non-telecentric DHM systems should compensate for the spherical phase factor recognized in
DHM and associated with a non-telecentric imaging system. Finally, the DHM configuration can operate in an image-plane configuration, meaning that in-focus DHM holograms
are recorded, so there is no need to apply numerical propagations to focus the amplitude and
phase images. However, if out-of-focus holograms are recorded, the user should numerically
propagate the complex amplitude distribution to provide in-focus images. Among the different
numerical propagation algorithms to reconstruct DHM images, the most used computational
approaches in DHM are the angular spectrum for short propagation distances [31, 32] and the
Fresnel transform algorithm for large propagation distances [33].
Still, each research group within the DHM community develops and implements its own
computational reconstruction algorithms based on its experimental digital holography (DH)
and DHM systems. Nonetheless, some research groups have developed and implemented
libraries and plugins to address the need for an open-source reconstruction toolbox in DHM
and DH. The main difference between DHM and DH systems is the utilization of a micro-
scopic imaging system in the object wave. 

\subsubsection{Background: Reconstruction algorithms for different optical
DHM configurations}

DHM systems are based on optical interferometry, which involves the digital recording of the
interference pattern (e.g., digital hologram) between the complex amplitude distribution scat-
tered by a microscopic object (e.g., the object wavefront) and a uniform reference wavefront.
DHM systems are colloquially known as hybrid imaging systems since the final reconstructed
amplitude or phase images are obtained after further computational processing of the digital
hologram. This second stage, known as the reconstruction stage, depends on the optical con-
figuration of the DHM system. This section thoroughly describes different DHM configura-
tions and their corresponding numerical reconstruction algorithms.
For simplicity, WE? have chosen a traditional DHM setup based on a Mach-Zehnder interfer-
ometer operating in transmission mode, see Fig 1(A) METTERE FIGURA. The light generated from a laser of wave-length λ is collimated by a converging lens (CL) and the resultant plane wavefront impinges a
cubic beam splitter (BS1) to generate the object (O) and the reference (R) wavefronts. A micro-
scopic imaging system is inserted in the object arm of the DHM system. This microscopic
imaging system comprises an infinity-corrected microscopic objective (MO) lens and a tube
lens (TL) which forms an image of the wavefront scattered by the microscopic sample with
amplitude distribution o(x,y). If the object is set at the front-focal plane (FFP) of the MO lens,
the image of the sample whose amplitude distribution is uIP (x,y) is then obtained at the back-
focal plane (BFP) of the TL. This plane is known as the microscope’s image plane (IP) [mettere e discutere eventualemente la formula (1)]. #fare il cammino ottico meglio

The hologram, captured at any plane from the IP, is the result of the interference between
the complex wavefield produced by the microscope at a distance zfrom the IP,....
discutere la formula (2) e (3) 
 The irradiance distribution of the hologram h(x,y;z) is ... (4)

On the other hand, the third
and fourth terms in Eq (4) encode the complex amplitude information of the object wavefront
scattered by the sample. These terms represent the object’s real and twin images. From Eq (4),
it is clear that the object information o(x,y), encoded in u(x,y;z) is mixed with other undesired
terms. The reconstruction algorithms aim to separate these undesired terms from the object
information to provide well-contrast amplitude and phase images from the complex in-focus
amplitude distribution uIP (x,y) = u(x,y;z= 0) with minimum distortions. Assuming that the
reference wavefront is plane, the 2D Fourier transform of the hologram H(u,v;z) is ....